\documentclass[polish, 8pt]{article}
\usepackage[top=2cm, bottom=2cm, left=3cm, right=3cm]{geometry}
\usepackage{graphicx} % Required for inserting images

\usepackage[polish]{babel}
\usepackage{polski}
\usepackage{placeins}
\usepackage{dirtytalk}
% Set page size and margins
% Replace `letterpaper' with `a4paper' for UK/EU standard size
\usepackage{csquotes}
% Useful packages
\usepackage{float}
\usepackage{amsmath}
\usepackage{graphicx}
\usepackage{hyperref}
\usepackage{graphicx}
\usepackage{listings}
\usepackage{xcolor}
\usepackage{amssymb}
\usepackage{enumitem}
\usepackage{float}
\definecolor{codegreen}{rgb}{0,0.6,0}
\definecolor{codegray}{rgb}{0.5,0.5,0.5}
\definecolor{codepurple}{rgb}{0.58,0,0.82}
\definecolor{backcolour}{rgb}{0.95,0.95,0.92}
\usepackage{amsmath, amssymb, tikz}
\lstdefinestyle{mystyle}{
    backgroundcolor=\color{backcolour},   
    commentstyle=\color{codegreen},
    keywordstyle=\color{magenta},
    numberstyle=\tiny\color{codegray},
    stringstyle=\color{codepurple},
    basicstyle=\ttfamily\footnotesize,
    breakatwhitespace=false,         
    breaklines=true,                 
    captionpos=b,                    
    keepspaces=true,                 
    numbers=left,                    
    numbersep=5pt,                  
    showspaces=false,                
    showstringspaces=false,
    showtabs=false,                  
    tabsize=2
}
\lstset{style=mystyle}
\hypersetup{
    colorlinks=true,
    linktoc=all     %set to all if you want both sections and subsections linked
}

\title{Rachunek Prawdopodobieństwa i Statystyka -- Kolokwium I}
\author{}
\date{}


\begin{document}
\maketitle

\section*{Zadanie 1}
Niech zmienna losowa dwuwymiarowa \((X, Y)\) ma rozkład ciągły o gęstości
\[
f(x, y) =
\begin{cases}
    c \, x^2 \, y, &\text{jeśli } x<0,\; y<1,\; x+y>0,\\
    0, &\text{w przeciwnym wypadku},
\end{cases}
\]
dla pewnej stałej \(c \in \mathbb{R}\).
\begin{enumerate}
    \item Wyznaczyć stałą \(c\).
    \item Wyznaczyć gęstości brzegowe: \(f_X(x)\) i \(f_Y(y)\).
    \item Wyznaczyć gęstości warunkowe: \(f(x\mid y)\) i \(f(y\mid x)\).
    \item Obliczyć prawdopodobieństwo
    \[
    P\!\bigl(X > -\tfrac{1}{4} \,\big\vert\, Y=\tfrac{1}{3}\bigr).
    \]
    \item Zbadać, czy zmienne \(X\) i \(Y\) są niezależne.
\end{enumerate}

\section*{Rozwiązanie}

\subsection*{1. Wyznaczenie stałej \(c\)}
Warunki na \((x,y)\) to:
\[
x < 0,\quad y < 1,\quad x + y > 0 \;\;\Longleftrightarrow\;\; y > -x.
\]
Ponadto, jeśli \(x \le -1\), to \(-x \ge 1\), co koliduje z warunkiem \(y < 1\). 
Dlatego faktycznie dziedzina \((x,y)\) zawęża się do
\[
-1 < x < 0, 
\quad -x \le y < 1.
\]
Funkcja gęstości musi się całkować do 1 po całej dziedzinie:
\[
\iint_{\mathcal{D}} f(x,y)\,dx\,dy 
=\iint_{\mathcal{D}} c\,x^2\,y \,dx\,dy = 1.
\]
Wygodniej całkować w kolejności \(dy\,dx\). Dla ustalonego \(x\in(-1,0)\) zmienna \(y\) przebiega \([\,-x,\,1]\). Zatem
\[
\int_{x=-1}^{0} \int_{y=-x}^{1} c \, x^2\,y \,dy\,dx = 1.
\]
Liczymy wartość tej podwójnej całki:
\[
\int_{-x}^{1} y \,dy 
= \left[ \frac{y^2}{2} \right]_{-x}^{1}
= \frac{1^2}{2} \;-\; \frac{(-x)^2}{2}
= \frac12 - \frac{x^2}{2}
= \frac12\bigl(1 - x^2\bigr).
\]
Zatem
\[
\int_{-1}^{0} c\,x^2 \left[\frac12\bigl(1 - x^2\bigr)\right] dx 
= c \cdot \frac12 \int_{-1}^{0} x^2 \bigl(1 - x^2\bigr)\,dx 
= c \cdot \frac12 \int_{-1}^{0} \bigl(x^2 - x^4\bigr)\,dx.
\]
Obliczamy kolejne całki:
\[
\int_{-1}^{0} x^2 \,dx = \left[\frac{x^3}{3}\right]_{-1}^{0} = 0 - \left(-\frac{1}{3}\right) = \frac{1}{3},
\]
\[
\int_{-1}^{0} x^4 \,dx = \left[\frac{x^5}{5}\right]_{-1}^{0} = 0 - \left(-\frac{1}{5}\right) = \frac{1}{5}.
\]
Wobec tego
\[
\int_{-1}^{0} \bigl(x^2 - x^4\bigr)\,dx 
= \frac{1}{3} - \frac{1}{5}
= \frac{5}{15} - \frac{3}{15}
= \frac{2}{15}.
\]
Zatem
\[
c \cdot \frac12 \cdot \frac{2}{15} = c \cdot \frac{1}{15} = 1
\quad\Longrightarrow\quad
c = 15.
\]
\[
\boxed{ c = 15. }
\]

\subsection*{2. Gęstości brzegowe \(f_X(x)\) oraz \(f_Y(y)\)}

\subsubsection*{Gęstość brzegowa \(f_X(x)\)}
Definicja:
\[
f_X(x) = \int_{-\infty}^{\infty} f(x,y)\,dy.
\]
Z uwagi na wsparcie gęstości, dla ustalonego \(x\in(-1,0)\) mamy \(y \in [-x,1]\). Wtedy:
\[
f_X(x) 
= \int_{y=-x}^{1} 15\,x^2\,y \,dy
= 15\,x^2 \int_{-x}^{1} y\,dy
\]
Jak wyżej wyliczyliśmy:
\[
\int_{-x}^{1} y\,dy 
= \frac12 \bigl(1 - x^2\bigr).
\]
Zatem
\[
f_X(x) = 15\,x^2 \cdot \frac12 \bigl(1 - x^2\bigr)
= \frac{15}{2}\,\bigl(x^2 - x^4\bigr),
\quad x\in(-1,0),
\]
a dla \(x\notin(-1,0)\) mamy \(f_X(x)=0\).  
Zapiszmy w formie zwarta:
\[
\boxed{
f_X(x) 
= \begin{cases}
\dfrac{15}{2}\,\bigl(x^2 - x^4\bigr), & -1 < x < 0,\\
0, & \text{w.p.w.}
\end{cases}
}
\]

\subsubsection*{Gęstość brzegowa \(f_Y(y)\)}
Definicja:
\[
f_Y(y) = \int_{-\infty}^{\infty} f(x,y)\,dx.
\]
Z warunków na dziedzinę:
\[
-1 < x < 0, \quad x + y > 0 \;\Longleftrightarrow\; x > -y,
\quad \text{oraz} \; y < 1.
\]
Aby móc znaleźć granice całkowania względem \(x\) dla ustalonego \(y\), zauważmy, że:
\[
x \text{ musi być w przedziale }(-1,0)\cap(-y,\infty).
\]
\emph{Kiedy} to nie jest pusty zbiór?
\begin{itemize}
\item Jeśli \(y < 0\), to \(-y > 0\). Wymaganie \(x > -y\) oznacza \(x> \text{(liczba dodatnia)}\), a to nie przecina się z warunkiem \(x<0\). Zatem dla \(y<0\) dziedzina jest pusta.
\item Jeśli \(0 \le y<1\), wówczas \(-y\le 0\). Przedział \((-y, \infty)\) przecięty z \((-1,0)\) daje \(\max(-1,-y)\), ale tu \(-y \le 0\). 
W rzeczywistości, przy \(0\le y<1\), warunek \(-1<x<0\) i \(x>-y\) sprowadza się do \(-y < x < 0\). 
\end{itemize}

Dodatkowo, w warunkach zadania jest \(y<1\). Dla \(y\ge 1\) nic nie wnosi (gęstość jest zerowa).  
Zatem faktycznie:
\[
0 < y < 1
\quad \Longrightarrow\quad
x \in (-y,\,0).
\]
Wtedy
\[
f_Y(y) 
= \int_{x=-y}^{0} 15\,x^2\,y \,dx 
= 15\,y \int_{-y}^{0} x^2\,dx.
\]
Obliczamy:
\[
\int_{-y}^{0} x^2\,dx 
= \left[\frac{x^3}{3}\right]_{-y}^{0}
= 0 - \left(-\frac{y^3}{3}\right)
= \frac{y^3}{3}.
\]
Zatem
\[
f_Y(y) = 15\,y \cdot \frac{y^3}{3}
= 5\,y^4,
\quad 0<y<1.
\]
Poza tym zakresem \(f_Y(y)=0\).  
Ostatecznie:
\[
\boxed{
f_Y(y) 
= \begin{cases}
5\,y^4, & 0<y<1, \\
0, & \text{w.p.w.}
\end{cases}
}
\]

\subsection*{3. Gęstości warunkowe}

\subsubsection*{Gęstość \(f_{X\mid Y}(x \mid y)\)}
Z definicji:
\[
f_{X\mid Y}(x \mid y) 
= \frac{f(x,y)}{f_Y(y)}.
\]
Dla \(y\in(0,1)\) oraz \(-y<x<0\) mamy \(f(x,y) = 15\,x^2\,y\), a \(f_Y(y)=5\,y^4\). Stąd
\[
f_{X\mid Y}(x\mid y)
= \frac{15\,x^2\,y}{5\,y^4}
= 3 \,\frac{x^2}{y^3},
\quad \text{dla } -y < x < 0.
\]
Poza tym przedziałem (w \(x\) lub \(y\)) mamy \(f_{X\mid Y}(x\mid y)=0\).  
Kontrola, czy się dobrze normalizuje (skrótowo):
\[
\int_{x=-y}^0 3\,\frac{x^2}{y^3}\,dx 
= \frac{3}{y^3} \int_{-y}^0 x^2\,dx 
= \frac{3}{y^3}\cdot \frac{y^3}{3} = 1.
\]

\subsubsection*{Gęstość \(f_{Y\mid X}(y \mid x)\)}
Z definicji:
\[
f_{Y\mid X}(y \mid x) 
= \frac{f(x,y)}{f_X(x)}.
\]
Dla \(x\in(-1,0)\) oraz \(-x<y<1\) mamy \(f(x,y) = 15\,x^2\,y\), a
\[
f_X(x) 
= \frac{15}{2}\,\bigl(x^2 - x^4\bigr).
\]
Zatem
\[
f_{Y\mid X}(y\mid x) 
= \frac{15\,x^2\,y}{\tfrac{15}{2}(x^2 - x^4)}
= \frac{2\,x^2\,y}{x^2 - x^4}
= 2\,y \,\frac{x^2}{x^2 (1 - x^2)}
= \frac{2\,y}{1 - x^2},
\quad \text{dla } -x<y<1.
\]
Poza tym zakresem jest 0.  
Można sprawdzić normalizację: 
\[
\int_{y=-x}^{1} \frac{2\,y}{1 - x^2}\,dy = \frac{2}{1 - x^2}\int_{-x}^{1} y\,dy = \ldots = 1.
\]

\subsection*{4. Obliczenie prawdopodobieństwa \(\displaystyle P\bigl(X>-\tfrac{1}{4}\mid Y=\tfrac{1}{3}\bigr)\)}

Korzystamy z gęstości warunkowej \(f_{X\mid Y}(x\mid y)\). Dla \(y=\tfrac13\) mamy:
\[
f_{X\mid Y}\!\Bigl(x \;\Big\vert\; y=\tfrac{1}{3}\Bigr) 
= 3\,\frac{x^2}{\bigl(\tfrac13\bigr)^3} 
= 3\,x^2 \cdot 27 
= 81\,x^2,
\quad \text{dla } -\tfrac13 < x < 0.
\]
(Sygnalizujemy, że dla \(Y=\tfrac13\), dziedzina \(X\) to przedział \(\bigl(-\tfrac13,0\bigr)\).)

Chcemy
\[
P\!\Bigl(X>-\tfrac{1}{4} \;\Big\vert\; Y=\tfrac{1}{3}\Bigr) 
= \int_{x=-\tfrac14}^{\,0} 81\,x^2 \,dx.
\]
Obliczamy:
\[
\int_{-1/4}^0 81\,x^2\,dx
= 81 \int_{-1/4}^0 x^2\,dx
= 81 \left[\frac{x^3}{3}\right]_{-1/4}^{0}
= 81 \Bigl(0 - \bigl(-\tfrac{1}{4}\bigr)^3/3\Bigr).
\]
Zauważmy, że \(\bigl(-\tfrac{1}{4}\bigr)^3 = -\tfrac{1}{64}\). Wobec tego:
\[
-\bigl(-\tfrac{1}{4}\bigr)^3 = \frac{1}{64}, 
\quad \text{więc} \quad 
-\bigl(-\tfrac{1}{4}\bigr)^3/3 = \frac{1}{64}\cdot \frac{1}{3} = \frac{1}{192}.
\]
Zatem 
\[
81 \cdot \frac{1}{192} 
= \frac{81}{192}
= \frac{81}{192} = \frac{27}{64}
\]
(po skróceniu przez 3).  

Otrzymujemy więc:
\[
\boxed{
P\!\Bigl(X > -\tfrac14 \;\Big\vert\; Y=\tfrac13\Bigr) 
= \frac{27}{64}.
}
\]

\subsection*{5. Niezależność zmiennych \(X\) i \(Y\)}
Zmienna losowa \(X\) i \(Y\) byłyby niezależne, gdyby zachodziło:
\[
f(x,y) \;=\; f_X(x)\,f_Y(y) 
\quad \text{dla każdej pary }(x,y).
\]
Tymczasem:
\[
f(x,y) = 15\,x^2\,y 
\quad\text{(w obszarze)}, 
\]
natomiast
\[
f_X(x)\,f_Y(y) 
= \Bigl[\tfrac{15}{2}\,(x^2 - x^4)\Bigr]
  \,\bigl[5\,y^4\bigr]
= \tfrac{75}{2}\,(x^2 - x^4)\,y^4.
\]
Te wyrażenia (w obszarze gdzie gęstość jest dodatnia) nie są sobie równe, 
co widać już po samym kształcie: jedno zależy od \(y\) w sposób liniowy (\(y\)) a drugie od \(y\) w sposób \(\,y^4\).  
Również sama dziedzina \(\{\,x+y>0,\dots\}\) to nie jest prostokąt (czyli nie dzieli się na niezależne przedziały \(\{x\in A\}\) i \(\{y\in B\}\)).  
\[
\boxed{
\text{Zmienne }X\text{ i }Y\text{ nie są niezależne.}
}
\]

\section*{Podsumowanie}
\begin{itemize}
\item Wyznaczona stała normalizacyjna: \(\displaystyle c = 15\).
\item Gęstości brzegowe:
\[
f_X(x) 
= 
\begin{cases}
\dfrac{15}{2}\,\bigl(x^2 - x^4\bigr), & -1<x<0,\\[6pt]
0, & \text{inaczej},
\end{cases}
\qquad
f_Y(y) 
= 
\begin{cases}
5\,y^4, & 0<y<1,\\
0, & \text{inaczej}.
\end{cases}
\]
\item Gęstości warunkowe:
\[
f_{X \mid Y}(x\mid y) = 
\begin{cases}
3\,\frac{x^2}{y^3}, & -y<x<0,\ 0<y<1,\\
0, & \text{inaczej},
\end{cases}
\quad
f_{Y \mid X}(y\mid x) = 
\begin{cases}
\dfrac{2\,y}{1 - x^2}, & -x<y<1,\ -1<x<0,\\
0, & \text{inaczej}.
\end{cases}
\]
\item Obliczone prawdopodobieństwo warunkowe:
\[
P\!\Bigl(X > -\tfrac14 \;\Big\vert\; Y=\tfrac13\Bigr) 
= \frac{27}{64}.
\]
\item Zmienna \(X\) i \(Y\) nie są niezależne.
\end{itemize}


\section*{Zadanie 2}

\section*{Zadanie 3}

\section*{Zadanie 4}
W zadaniu mamy do dyspozycji:
\[
n = 16,\quad
\overline{X} = 3000 \text{ (średnia z próby)},\quad
s = 20 \text{ (odchylenie standardowe z próby)},\quad
1 - \alpha = 0{,}90 \quad (\text{poziom ufności}).
\]
Z założenia, że rozkład czasu życia żarówek jest normalny, a odchylenie standardowe w populacji nie jest znane, do konstrukcji przedziału ufności dla średniej wykorzystujemy rozkład \emph{t-Studenta}.

\subsection*{Wyjaśnienie doboru statystyki}
Ponieważ nie znamy odchylenia standardowego w populacji, a jedynie jego estymator $s$ wyznaczony z próby, stosujemy statystykę:
\[
T = \frac{\overline{X} - \mu}{s / \sqrt{n}},
\]
która ma rozkład \emph{t-Studenta} z $n-1$ stopniami swobody (tj. w naszym przypadku $df = 16 - 1 = 15$). 
W przypadku, gdybyśmy znali odchylenie standardowe w populacji (oznaczane zwykle $\sigma$), wtedy moglibyśmy wykorzystać rozkład normalny (statystykę $Z$). 
Tutaj jednak musimy użyć rozkładu $t$, ponieważ $\sigma$ jest nieznane, a $n$ jest stosunkowo małe.

\subsection*{Wzór na przedział ufności}
Przedział ufności $100(1-\alpha)\%$ dla średniej $\mu$ w przypadku nieznanego $\sigma$ i małej próby ma postać:
\[
\left(\,\overline{X} - t_{\alpha/2,\,n-1}\,\frac{s}{\sqrt{n}},\;\;
       \overline{X} + t_{\alpha/2,\,n-1}\,\frac{s}{\sqrt{n}} \right),
\]
gdzie $t_{\alpha/2,\,n-1}$ to \emph{kwantyl} (wartość krytyczna) rozkładu t-Studenta o $n-1$ stopniach swobody, wyznaczony tak, aby w obu ogonach łącznie było $\alpha$, czyli w jednym ogonie jest $\alpha/2$. 

\subsection*{Obliczenie wartości krytycznej}
Dla poziomu ufności $1-\alpha = 0{,}90$, otrzymujemy $\alpha = 0{,}10$, a więc $\alpha/2 = 0{,}05$. 
Szukamy zatem kwantyla $t_{0,05, 15}$, odpowiadającego $15$ stopniom swobody i granicy $0,95$ rozkładu. 
Z tablic rozkładu t-Studenta otrzymujemy w przybliżeniu:
\[
t_{\alpha/2,\,15} = t_{0,05,\,15} \approx 1{,}753.
\]

\subsection*{Obliczenie szerokości przedziału}
\[
\text{Błąd standardowy} = \frac{s}{\sqrt{n}} = \frac{20}{\sqrt{16}} = \frac{20}{4} = 5.
\]
\[
\text{Szerokość półprzedziału ufności} 
= t_{\alpha/2,\,15}\,\frac{s}{\sqrt{n}}
= 1{,}753 \cdot 5
\approx 8{,}765.
\]

\subsection*{Konstrukcja przedziału ufności}
\[
\left(\,\overline{X} - t_{\alpha/2,\,15}\,\frac{s}{\sqrt{n}},
\quad \overline{X} + t_{\alpha/2,\,15}\,\frac{s}{\sqrt{n}}\right) 
= \left(3000 - 8{,}765,\; 3000 + 8{,}765\right).
\]
\[
\approx (2991{,}235,\; 3008{,}765).
\]

Zatem przedział ufności $90\%$ dla wartości średniej czasu życia żarówki to:
\[
\boxed{
(2991{,}24,\; 3008{,}76).
}
\]
\section*{Zadanie 5}

\end{document}

